% ============================================================
% Come si usano, argomento 23: stima lineare MMSE e gradiente
% ============================================================
\argomento{Stima lineare MMSE e gradiente}{Formule a p.~\pageref*{f:lmmse}}

\begin{scelta}{Come scegliere la formula}
\renewcommand{\arraystretch}{1.25}
\begin{tabular}{@{}P{0.55\linewidth}P{0.39\linewidth}@{}}
  \toprule
  \textbf{Se il testo\ldots} & \textbf{si usa\ldots} \\
  \midrule
  chiede il miglior stimatore lineare, dando medie e covarianze & \vv{lmmse} \\
  misura un parametro con uno strumento rumoroso & \vv{misurarumorosa} \\
  chiede di trovare i coefficienti per iterazioni, o il passo che fa convergere & \vv{gradientealg}, con gli
    autovalori \vv{autovalori} \\
  \bottomrule
\end{tabular}

Servono solo medie e covarianze: la matrice $\matr{M}$ dei dati, il vettore $\vett{s}$ delle covarianze fra dati e
parametro, la media e la varianza del parametro. Conviene controllare le dimensioni: $\vett{s}^T \matr{M}^{-1}
\vett{s}$ è un numero.
\end{scelta}

\scheda{lmmse}{
  \usoquando{Si vuole stimare un parametro aleatorio con una combinazione lineare dei dati, conoscendo solo medie e
    covarianze (non tutta la densità a posteriori).}
  \usocome{Si costruiscono $\matr{M}$ e $\vett{s}$, si risolve $\matr{M}\vett{a} = \vett{s}$ e si sceglie $b$ in modo
    che la stima abbia la media giusta.}
  \usoesempio{$\Theta$ con media 1 e varianza 1, due misure $X_i = \Theta + N_i$ con rumori indipendenti di varianza
    1: $\matr{M} = \begin{bmatrix} 2 & 1 \\ 1 & 2 \end{bmatrix}$, $\vett{s} = [1, 1]^T$, $\vett{a} = \left[
    \frac{1}{3}, \frac{1}{3} \right]^T$, $b = \frac{1}{3}$. Quindi $\hat\Theta = \frac{X_1 + X_2 + 1}{3}$, con errore
    $1 - \frac{2}{3} = \frac{1}{3}$.}
}

\scheda{misurarumorosa}{
  \usoquando{Un parametro con media e varianza note a priori viene misurato una volta, con un rumore additivo.}
  \usocome{La stima è una media pesata fra la misura e la media a priori: con poco rumore ci si fida della misura, con
    tanto rumore di ciò che si sapeva prima.}
  \usoesempio{$\mu_\Theta = 2$, $\sigma_\Theta = 1$, $\sigma_N = 0{,}7$: $a = \frac{1}{1{,}49} \simeq 0{,}67$; con la
    misura $X = 3$ la stima è $0{,}67 \cdot 3 + 0{,}33 \cdot 2 \simeq 2{,}67$. L'errore scende da $0{,}49$ a
    $0{,}33$.}
}

\scheda{gradientealg}{
  \usoquando{Invertire $\matr{M}$ costa troppo, o le statistiche cambiano nel tempo: si scende verso il minimo a piccoli
    passi.}
  \usocome{Si calcolano gli autovalori di $\matr{M}$: il passo deve stare sotto $\frac{2}{\lambda_{\text{MAX}}}$, e con
    autovalori molto diversi la convergenza è lenta lungo quello piccolo.}
  \usoesempio{$\matr{M} = \begin{bmatrix} 2 & 0{,}9 \\ 0{,}9 & 1 \end{bmatrix}$: autovalori $2{,}53$ e $0{,}47$, quindi
    serve $\gamma < 0{,}79$. Con $\gamma = 0{,}25$ i fattori sono $0{,}37$ e $0{,}88$ (lento); con $\gamma = 0{,}72$ il
    primo è $-0{,}82$ e l'algoritmo procede a zig-zag.}
}
